\section{Spectral spaces}
\label{section-spectral}

\noindent
The material in this section is taken from \cite{Hochster} and
\cite{Hochster-thesis}. In his thesis Hochster proves (among
other things) that the spectral spaces are exactly the topological
spaces that occur as the spectrum of a ring.

\begin{definition}
\label{definition-spectral-space}
A topological space $X$ is called {\it spectral} if it is sober,
quasi-compact, the intersection of two quasi-compact opens is
quasi-compact, and the collection of quasi-compact opens forms a
basis for the topology. A continuous map $f : X \to Y$ of spectral
spaces is called {\it spectral} if the inverse image of a quasi-compact
open is quasi-compact.
\end{definition}

\noindent
In other words a continuous map of spectral spaces is spectral if and
only if it is quasi-compact (Definition \ref{definition-quasi-compact}).

\medskip\noindent
Let $X$ be a spectral space. The {\it constructible topology} on $X$
is the topology which has as a subbase of opens the sets $U$ and $U^c$
where $U$ is a quasi-compact open of $X$. Note that since $X$ is spectral
an open $U \subset X$ is retrocompact if and only if $U$ is quasi-compact.
Hence the constructible topology can also be characterized as the coarsest
topology such that every constructible subset of $X$ is both open and closed
(see Section \ref{section-constructible} for definitions and properties of
constructible sets).
It follows that a subset of $X$ is open, resp.\ closed in the constructible
topology if and only if it is a union, resp.\ intersection of constructible
subsets. Since the collection of quasi-compact opens is a basis for the topology
on $X$ we see that the constructible topology is stronger than the
given topology on $X$.

\begin{lemma}
\label{lemma-constructible-hausdorff-quasi-compact}
Let $X$ be a spectral space. The constructible topology is
Hausdorff, totally disconnected, and quasi-compact.
\end{lemma}

\begin{proof}
Let $x, y \in X$ with $x \not = y$. Since $X$ is sober, there
is an open subset $U$ containing exactly one of the two points $x, y$.
Say $x \in U$. We may replace $U$ by a quasi-compact open
neighbourhood of $x$ contained in $U$. Then $U$ and $U^c$ are open and
closed in the constructible topology. Hence $X$ is Hausdorff in the
constructible topology because $x \in U$ and $y \in U^c$ are
disjoint opens in the constructible topology. The existence of $U$
also implies $x$ and $y$ are in distinct connected components in the
constructible topology, whence $X$ is totally disconnected in the
constructible topology.

\medskip\noindent
Let $\mathcal{B}$ be the collection of subsets $B \subset X$
with $B$ either quasi-compact open or closed with quasi-compact
complement. If $B \in \mathcal{B}$ then $B^c \in \mathcal{B}$.
It suffices to show every covering $X = \bigcup_{i \in I} B_i$
with $B_i \in \mathcal{B}$ has a finite refinement, see
Lemma \ref{lemma-subbase-theorem}.
Taking complements we see that we have to show that any family
$\{B_i\}_{i \in I}$ of elements of $\mathcal{B}$
such that $B_{i_1} \cap \ldots \cap B_{i_n} \not = \emptyset$
for all $n$ and all $i_1, \ldots, i_n \in I$
has a common point of intersection. We may and do assume
$B_i \not = B_{i'}$ for $i \not = i'$.

\medskip\noindent
To get a contradiction assume $\{B_i\}_{i \in I}$ is a family
of elements of $\mathcal{B}$ having the finite intersection property
but empty intersection. An application of Zorn's lemma shows that we
may assume our family is maximal (details omitted).
Let $I' \subset I$ be those indices such that
$B_i$ is closed and set $Z = \bigcap_{i \in I'} B_i$.
This is a closed subset of $X$ which is nonempty by Lemma
\ref{lemma-intersection-closed-in-quasi-compact}.
If $Z$ is reducible, then we can write $Z = Z' \cup Z''$
as a union of two closed subsets, neither equal to $Z$. This means in
particular that we can find a quasi-compact open $U' \subset X$ meeting
$Z'$ but not $Z''$. Similarly, we can find a quasi-compact open
$U'' \subset X$ meeting $Z''$ but not $Z'$. Set $B' = X \setminus U'$ and
$B'' = X \setminus U''$. Note that $Z'' \subset B'$ and $Z' \subset B''$.
If there exist a finite number of indices $i_1, \ldots, i_n \in I$ such
that $B' \cap B_{i_1} \cap \ldots \cap B_{i_n} = \emptyset$
as well as a finite number of indices $j_1, \ldots, j_m \in I$ such that
$B'' \cap B_{j_1} \cap \ldots \cap B_{j_m} = \emptyset$
then we find that
$Z \cap B_{i_1} \cap \ldots \cap B_{i_n} \cap B_{j_1} \cap \ldots \cap B_{j_m}
= \emptyset$.
However, the set
$B_{i_1} \cap \ldots \cap B_{i_n} \cap B_{j_1} \cap \ldots \cap B_{j_m}$
is quasi-compact hence we would find a finite number of
indices $i'_1, \ldots, i'_l \in I'$ with
$B_{i_1} \cap \ldots \cap B_{i_n} \cap B_{j_1} \cap \ldots \cap
B_{j_m} \cap B_{i'_1} \cap \ldots \cap B_{i'_l} = \emptyset$, a contradiction.
Thus we see that we may add either $B'$ or $B''$ to the given family
contradicting maximality. We conclude that $Z$ is irreducible. However,
this leads to a contradiction as well, as now every nonempty (by the
same argument as above) open $Z \cap B_i$ for $i \in I \setminus I'$
contains the unique generic point of $Z$. This contradiction proves the lemma.
\end{proof}

\begin{lemma}
\label{lemma-fibres-spectral-map-quasi-compact}
Let $f : X \to Y$ be a spectral map of spectral spaces. Then
\begin{enumerate}
\item $f$ is continuous in the constructible topology,
\item the fibres of $f$ are quasi-compact, and
\item the image is closed in the constructible topology.
\end{enumerate}
\end{lemma}

\begin{proof}
Let $X'$ and $Y'$ denote $X$ and $Y$ endowed with the constructible topology
which are quasi-compact Hausdorff spaces by
Lemma \ref{lemma-constructible-hausdorff-quasi-compact}.
Part (1) says $X' \to Y'$ is continuous and follows immediately from the
definitions. Part (3) follows as $f(X')$ is a quasi-compact subset of the
Hausdorff space $Y'$, see Lemma \ref{lemma-quasi-compact-in-Hausdorff}.
We have a commutative diagram
$$
\xymatrix{
X' \ar[r] \ar[d] & X \ar[d] \\
Y' \ar[r] & Y
}
$$
of continuous maps of topological spaces. Since $Y'$ is Hausdorff
we see that the fibres $X'_y$ are closed in $X'$. As $X'$ is quasi-compact
we see that $X'_y$ is quasi-compact
(Lemma \ref{lemma-closed-in-quasi-compact}).
As $X'_y \to X_y$ is a surjective continuous map we conclude
that $X_y$ is quasi-compact (Lemma \ref{lemma-image-quasi-compact}).
\end{proof}

\begin{lemma}
\label{lemma-spectral-if-continuous-wrt-constructible-top}
Let $X$ and $Y$ be spectral spaces. Let $f : X \to Y$ be a continuous map.
Then $f$ is spectral if and only if $f$ is continuous in the constructible
topology.
\end{lemma}

\begin{proof}
The only if part of this is
Lemma \ref{lemma-fibres-spectral-map-quasi-compact}.
Assume $f$ is continuous in the constructible topology.
Let $V \subset Y$ be quasi-compact open.
Then $V$ is open and closed in the constructible topology.
Hence $f^{-1}(V)$ is open and closed in the constructible topology.
Hence $f^{-1}(V)$ is quasi-compact in the constructible topology
as $X$ is quasi-compact in the constructible topology by
Lemma \ref{lemma-constructible-hausdorff-quasi-compact}.
Since the identity $f^{-1}(V) \to f^{-1}(V)$ is surjective
and continuous from the constructible topology to the usual
topology, we conclude that $f^{-1}(V)$ is quasi-compact
in the topology of $X$ by Lemma \ref{lemma-image-quasi-compact}.
This finishes the proof.
\end{proof}

\begin{lemma}
\label{lemma-spectral-sub}
Let $X$ be a spectral space. Let $E \subset X$ be closed in the constructible
topology (for example constructible or closed). Then $E$ with the induced
topology is a spectral space.
\end{lemma}

\begin{proof}
Let $Z \subset E$ be a closed irreducible subset. Let $\eta$ be the generic
point of the closure $\overline{Z}$ of $Z$ in $X$. To prove that $E$
is sober, we show that $\eta \in E$. If not, then since $E$ is closed
in the constructible topology, there exists a constructible subset
$F \subset X$ such that $\eta \in F$ and $F \cap E = \emptyset$.
By Lemma \ref{lemma-generic-point-in-constructible} this implies
$F \cap \overline{Z}$ contains a nonempty open subset of $\overline{Z}$.
But this is impossible as $\overline{Z}$ is the closure of $Z$ and
$Z \cap F = \emptyset$.

\medskip\noindent
Since $E$ is closed in the constructible topology, it is quasi-compact
in the constructible topology
(Lemmas \ref{lemma-closed-in-quasi-compact} and
\ref{lemma-constructible-hausdorff-quasi-compact}). Hence a fortiori it is
quasi-compact in the topology coming from $X$. If $U \subset X$
is a quasi-compact open, then $E \cap U$ is closed in the constructible
topology, hence quasi-compact (as seen above). It follows that the
quasi-compact open subsets of $E$ are the intersections $E \cap U$
with $U$ quasi-compact open in $X$. These form a basis for the topology.
Finally, given two $U, U' \subset X$ quasi-compact opens, the intersection
$(E \cap U) \cap (E \cap U') = E \cap (U \cap U')$ and $U \cap U'$
is quasi-compact as $X$ is spectral. This finishes the proof.
\end{proof}

\begin{lemma}
\label{lemma-constructible-stable-specialization-closed}
Let $X$ be a spectral space. Let $E \subset X$ be a subset closed
in the constructible topology (for example constructible).
\begin{enumerate}
\item If $x \in \overline{E}$, then $x$ is the specialization of a point of
$E$.
\item If $E$ is stable under specialization, then $E$ is closed.
\item If $E' \subset X$ is open in the constructible topology
(for example constructible) and stable under generalization, then $E'$ is open.
\end{enumerate}
\end{lemma}

\begin{proof}
Proof of (1). Let $x \in \overline{E}$. Let $\{U_i\}$ be the set of
quasi-compact open neighbourhoods of $x$. A finite intersection of the
$U_i$ is another one. The intersection $U_i \cap E$ is nonempty for all
$i$. Since the subsets $U_i \cap E$ are closed in the constructible topology
we see that $\bigcap (U_i \cap E)$ is nonempty by
Lemma \ref{lemma-constructible-hausdorff-quasi-compact} and
Lemma \ref{lemma-intersection-closed-in-quasi-compact}.
Since $\{U_i\}$ is a
fundamental system of open neighbourhoods of $x$, we see that
$\bigcap U_i$ is the set of generalizations of $x$. Thus
$x$ is a specialization of a point of $E$.

\medskip\noindent
Part (2) is immediate from (1).

\medskip\noindent
Proof of (3). Assume $E'$ is as in (3). The complement of $E'$ is closed
in the constructible topology (Lemma \ref{lemma-constructible})
and closed under specialization
(Lemma \ref{lemma-open-closed-specialization}).
Hence the complement is closed by (2), i.e., $E'$ is open.
\end{proof}

\begin{lemma}
\label{lemma-two-points}
Let $X$ be a spectral space. Let $x, y \in X$. Then either there exists
a third point specializing to both $x$ and $y$, or there exist disjoint
open neighbourhoods containing $x$ and $y$.
\end{lemma}

\begin{proof}
Let $\{U_i\}$ be the set of quasi-compact open neighbourhoods of $x$.
A finite intersection of the $U_i$ is another one.
Let $\{V_j\}$ be the set of quasi-compact open neighbourhoods of $y$.
A finite intersection of the $V_j$ is another one.
If $U_i \cap V_j$ is empty for some $i, j$ we are done.
If not, then the intersection $U_i \cap V_j$ is nonempty
for all $i$ and $j$. The sets $U_i \cap V_j$ are closed in the constructible
topology on $X$. By
Lemma \ref{lemma-constructible-hausdorff-quasi-compact}
we see that $\bigcap (U_i \cap V_j)$ is nonempty
(Lemma \ref{lemma-intersection-closed-in-quasi-compact}).
Since $X$ is a sober space and $\{U_i\}$ is a
fundamental system of open neighbourhoods of $x$, we see that
$\bigcap U_i$ is the set of generalizations of $x$.
Similarly, $\bigcap V_j$ is the set of generalizations of $y$.
Thus any element of $\bigcap (U_i \cap V_j)$ specializes to both
$x$ and $y$.
\end{proof}

\begin{lemma}
\label{lemma-characterize-profinite-spectral}
Let $X$ be a spectral space. The following are equivalent:
\begin{enumerate}
\item $X$ is profinite,
\item $X$ is Hausdorff,
\item $X$ is totally disconnected,
\item every quasi-compact open is closed,
\item there are no nontrivial specializations between points,
\item every point of $X$ is closed,
\item every point of $X$ is the generic point of an irreducible component
of $X$,
\item the constructible topology equals the given topology on $X$, and
\item add more here.
\end{enumerate}
\end{lemma}

\begin{proof}
Lemma \ref{lemma-profinite} shows the implication (1) $\Rightarrow$ (3).
Irreducible components are closed, so if $X$ is totally disconnected, then
every point is closed. So (3) implies (6). The equivalence of (6) and (5)
is immediate, and (6) $\Leftrightarrow$ (7) holds because $X$ is sober.
Assume (5). Then all constructible subsets of $X$ are closed
(Lemma \ref{lemma-constructible-stable-specialization-closed}), in
particular all quasi-compact opens are closed. So (5) implies (4).
Since $X$ is sober, for any two points there is a quasi-compact open
containing exactly one of them, hence (4) implies (2).
Parts (4) and (8) are equivalent by the definition of the
constructible topology.
It remains to prove (2) implies (1). Suppose $X$ is Hausdorff. Every
quasi-compact open is also closed
(Lemma \ref{lemma-quasi-compact-in-Hausdorff}). This implies $X$
is totally disconnected. Hence it is profinite, by
Lemma \ref{lemma-profinite}.
\end{proof}

\begin{lemma}
\label{lemma-spectral-pi0}
If $X$ is a spectral space, then $\pi_0(X)$ is a profinite space.
\end{lemma}

\begin{proof}
Combine Lemmas \ref{lemma-connected-component-intersection} and
\ref{lemma-pi0-profinite}.
\end{proof}

\begin{lemma}
\label{lemma-product-spectral-spaces}
The product of two spectral spaces is spectral.
\end{lemma}

\begin{proof}
Let $X$, $Y$ be spectral spaces. Denote $p : X \times Y \to X$ and
$q : X \times Y \to Y$ the projections. Let $Z \subset X \times Y$ be a
closed irreducible subset. Then $p(Z) \subset X$ is irreducible
and $q(Z) \subset Y$ is irreducible. Let $x \in X$ be the generic
point of the closure of $p(X)$ and let $y \in Y$ be the generic
point of the closure of $q(Y)$. If $(x, y) \not \in Z$, then
there exist opens $x \in U \subset X$, $y \in V \subset Y$ such
that $Z \cap U \times V = \emptyset$. Hence $Z$ is contained
in $(X \setminus U) \times Y \cup X \times (Y \setminus V)$.
Since $Z$ is irreducible, we see that either
$Z \subset (X \setminus U) \times Y$ or $Z \subset X \times (Y \setminus V)$.
In the first case $p(Z) \subset (X \setminus U)$ and in the
second case $q(Z) \subset (Y \setminus V)$. Both cases are absurd
as $x$ is in the closure of $p(Z)$ and $y$ is in the closure of
$q(Z)$. Thus we conclude that $(x, y) \in Z$, which means that
$(x, y)$ is the generic point for $Z$.

\medskip\noindent
A basis of the topology of $X \times Y$ are the opens of the form
$U \times V$ with $U \subset X$ and $V \subset Y$ quasi-compact open
(here we use that $X$ and $Y$ are spectral). Then $U \times V$ is
quasi-compact as the product of quasi-compact spaces is quasi-compact.
Moreover, any quasi-compact open of $X \times Y$ is a finite union
of such quasi-compact rectangles $U \times V$. It follows that
the intersection of two such is again quasi-compact
(since $X$ and $Y$ are spectral). This concludes the proof.
\end{proof}

\begin{lemma}
\label{lemma-spectral-bijective}
Let $f : X \to Y$ be a continuous map of topological spaces. If
\begin{enumerate}
\item $X$ and $Y$ are spectral,
\item $f$ is spectral and bijective, and
\item generalizations (resp.\ specializations) lift along $f$.
\end{enumerate}
Then $f$ is a homeomorphism.
\end{lemma}

\begin{proof}
Since $f$ is spectral it defines a continuous map between $X$ and $Y$ in
the constructible topology. By
Lemmas \ref{lemma-constructible-hausdorff-quasi-compact} and
\ref{lemma-bijective-map}
it follows that $X \to Y$ is a homeomorphism in the constructible
topology. Let $U \subset X$ be quasi-compact open.
Then $f(U)$ is constructible in $Y$. Let $y \in Y$ specialize
to a point in $f(U)$. By the last assumption we see that $f^{-1}(y)$
specializes to a point of $U$. Hence $f^{-1}(y) \in U$. Thus $y \in f(U)$.
It follows that $f(U)$ is open, see
Lemma \ref{lemma-constructible-stable-specialization-closed}.
Whence $f$ is a homeomorphism.
To prove the lemma in case specializations lift along $f$
one shows instead that $f(Z)$ is closed if $X \setminus Z$ is a
quasi-compact open of $X$.
\end{proof}

\begin{lemma}
\label{lemma-directed-inverse-limit-finite-sober-spectral-spaces}
The inverse limit of a directed inverse system of finite sober
topological spaces is a spectral topological space.
\end{lemma}

\begin{proof}
Let $I$ be a directed set. Let $X_i$ be an inverse
system of finite sober spaces over $I$. Let $X = \lim X_i$ which exists
by Lemma \ref{lemma-limits}. As a set $X = \lim X_i$. Denote
$p_i : X \to X_i$ the projection.
Because $I$ is directed we may apply Lemma \ref{lemma-describe-limits}.
A basis for the topology is given by the opens
$p_i^{-1}(U_i)$ for $U_i \subset X_i$ open. Since an open covering of
$p_i^{-1}(U_i)$ is in particular an open covering in the profinite
topology, we conclude that $p_i^{-1}(U_i)$ is quasi-compact.
Given $U_i \subset X_i$ and $U_j \subset X_j$, then
$p_i^{-1}(U_i) \cap p_j^{-1}(U_j) = p_k^{-1}(U_k)$
for some $k \geq i, j$ and open $U_k \subset X_k$. Finally, if $Z \subset X$
is irreducible and closed, then $p_i(Z) \subset X_i$ is irreducible
and therefore has a unique generic point $\xi_i$ (because $X_i$
is a finite sober topological space). Then $\xi = \lim \xi_i$ is a
generic point of $Z$ (it is a point of $Z$ as $Z$ is closed).
This finishes the proof.
\end{proof}

\begin{lemma}
\label{lemma-spectral-closed-in-product-two-point-space}
Let $W$ be the topological space with two points, one closed,
the other not. A topological space is spectral if and only if
it is homeomorphic to a subspace of a product of
copies of $W$ which is closed in the constructible topology.
\end{lemma}

\begin{proof}
Write $W = \{0, 1\}$ where $0$ is a specialization of $1$ but not vice versa.
Let $I$ be a set. The space $\prod_{i \in I} W$ is spectral by
Lemma \ref{lemma-directed-inverse-limit-finite-sober-spectral-spaces}.
Thus we see that a subspace of $\prod_{i \in I} W$ closed in the
constructible topology is
a spectral space by Lemma \ref{lemma-spectral-sub}.

\medskip\noindent
For the converse, let $X$ be a spectral space. Let $U \subset X$ be a
quasi-compact open. Consider the continuous map
$$
f_U : X \longrightarrow W
$$
which maps every point in $U$ to $1$ and every point in $X \setminus U$ to $0$.
Taking the product of these maps we obtain a continuous map
$$
f = \prod f_U : X \longrightarrow \prod\nolimits_U W
$$
By construction the map $f : X \to Y$ is spectral. By
Lemma \ref{lemma-fibres-spectral-map-quasi-compact}
the image of $f$ is closed in the constructible topology.
If $x', x \in X$ are distinct, then since $X$ is sober either $x'$
is not a specialization of $x$ or conversely. In either case (as the
quasi-compact opens form a basis for the topology of $X$) there
exists a quasi-compact open $U \subset X$ such that $f_U(x') \not = f_U(x)$.
Thus $f$ is injective. Let $Y = f(X)$ endowed with the induced topology.
Let $y' \leadsto y$ be a specialization in $Y$ and say
$f(x') = y'$ and $f(x) = y$. Arguing as above we see that
$x' \leadsto x$, since otherwise there is a $U$ such that
$x \in U$ and $x' \not \in U$, which would imply
$f_U(x') \not \leadsto f_U(x)$.
We conclude that $f : X \to Y$ is a homeomorphism by
Lemma \ref{lemma-spectral-bijective}.
\end{proof}

\begin{lemma}
\label{lemma-spectral-inverse-limit-finite-sober-spaces}
A topological space is spectral if and only if it is a directed
inverse limit of finite sober topological spaces.
\end{lemma}

\begin{proof}
One direction is given by
Lemma \ref{lemma-directed-inverse-limit-finite-sober-spectral-spaces}.
For the converse, assume $X$ is spectral. Then we
may assume $X \subset \prod_{i \in I} W$ is a subset closed
in the constructible topology where $W = \{0, 1\}$ as in
Lemma \ref{lemma-spectral-closed-in-product-two-point-space}.
We can write
$$
\prod\nolimits_{i \in I} W =
\lim_{J \subset I\text{ finite }} \prod\nolimits_{j \in J} W
$$
as a cofiltered limit.
For each $J$, let $X_J \subset \prod_{j \in J} W$ be the image of $X$.
Then we see that $X = \lim X_J$ as sets because $X$ is closed in
the product with the constructible topology (detail omitted).
A formal argument (omitted) on limits shows that $X = \lim X_J$ as topological
spaces.
\end{proof}

\begin{lemma}
\label{lemma-Noetherian-goes-to-spectral}
Let $X$ be a topological space and let $c : X \to X'$ be the universal
map from $X$ to a sober topological space, see
Lemma \ref{lemma-make-sober}.
\begin{enumerate}
\item If $X$ is quasi-compact, so is $X'$.
\item If $X$ is quasi-compact, has a basis of quasi-compact opens,
and the intersection of two quasi-compact opens is quasi-compact, then
$X'$ is spectral.
\item If $X$ is Noetherian, then $X'$ is a Noetherian spectral space.
\end{enumerate}
\end{lemma}

\begin{proof}
Let $U \subset X$ be open and let $U' \subset X'$ be the corresponding
open, i.e., the open such that $c^{-1}(U') = U$.
Then $U$ is quasi-compact if and only if $U'$ is quasi-compact,
as pulling back by $c$ is a bijection between the opens of
$X$ and $X'$ which commutes with unions. This in particular proves (1).

\medskip\noindent
Proof of (2). It follows from the above that $X'$ has a basis of
quasi-compact opens. Since $c^{-1}$ also
commutes with intersections of pairs of opens, we see that
the intersection of two quasi-compact opens $X'$ is quasi-compact.
Finally, $X'$ is quasi-compact by (1) and sober by construction.
Hence $X'$ is spectral.

\medskip\noindent
Proof of (3). It is immediate that $X'$ is Noetherian as this is defined
in terms of the acc for open subsets which holds for $X$. We have
already seen in (2) that $X'$ is spectral.
\end{proof}




